\section{Limitations and Conclusion}
\label{sec:conclusion}
RAI's checkpoints grow with the unresolved forks they keep, and locks, late supports, and overlap witnesses must be kept while a lock they can release survives; we give no garbage-collection policy beyond the prototype's four-epoch cap. The committee model is authenticated and equal-weight with an externally kept fault bound. The evaluation runs on one host, so the delays it measures are dominated by CPU contention, each cell ran three times, and process kills test crash consistency, not power loss. The models cover at most two positions and three epochs, and the termination rules only $N=4$.

Within these limits RAI does what it set out to do. It keeps account-local payments safe across committee changes without stopping: the next epoch votes while the previous one closes, immutable reports preserve whatever the old epoch could still confirm, and four overlap rules keep the two epochs from contradicting each other. It keeps at most one side of every contested position, identically everywhere, and finalizes that side within two epochs of normal operation, so the owner always gets the account back. Against a stop-and-drain transition built from the same code, it keeps confirming through every boundary and shortens the latency tail severalfold, most of all when a faulty validator leads the closes. Against a system without epochs, at the same quorum, it finalizes the forks that system leaves open and keeps confirming at almost three times its rate when forks are frequent, but without forks its closes compete with voting for the cores of a shared host, so it confirms 5\% fewer blocks per second and has a longer 95th percentile. In return it can change its validators, and its protocol: every validator agrees on the epoch boundaries and the checkpoint that closes each epoch, so a protocol change can take effect at an agreed epoch, where Nano appends an epoch block to every account.
